\section{Introduction}\label{sec:introduction}

Equatorial Plasma Bubbles (EPBs) are large regions of depleted plasma density that form in the equatorial F region after sunset through the generalized Rayleigh-Taylor instability. Once the Sun sets, recombination quickly erodes the bottomside ionosphere and leaves a steep upward density gradient beneath the F layer. When the pre-reversal enhancement of the eastward electric field is strong, this configuration becomes unstable: low-density plasma rises into the topside ionosphere as plumes that extend to low latitudes and drift eastward. The small-scale irregularities inside these depletions cause amplitude and phase scintillation in radio signals, degrading GNSS positioning and satellite communications \cite{zhong2025automatic,okoh2025bootstrapping,githio2024machine}.

From the ground, EPBs are usually observed with All-Sky Imagers (ASIs) fitted with narrow-band filters at $630.0\text{ nm}$, the atomic oxygen red line emitted at about $250\text{--}300\text{ km}$ \cite{okoh2025bootstrapping,zhong2025automatic}. Because the emission is weaker inside the depletions, bubbles appear as dark bands elongated in the north--south direction. At an observatory, finding them still means looking through hundreds of images per night. The task is slow, and it becomes ambiguous when faint depletions overlap with moonlight, horizon brightening from the Van Rhijn effect, the Milky Way, or passing clouds.

Several recent works have applied machine learning to this problem. Convolutional Neural Networks (CNNs) and segmentation models have been used to classify EPB occurrence, track medium-scale traveling ionospheric disturbances, and build multi-year climatologies \cite{srisamoodkham2022detecting,liu2022statistical,sedlak2023analysis,okoh2025bootstrapping,zhong2025automatic,zhong2025statistical}, and other studies combined CNN feature extractors with Support Vector Machines (SVMs) or with collocated GNSS scintillation data \cite{thanakulketsarat2023classification,thanakulketsarat2023equatorial,githio2024machine}. Deep networks, however, are costly to train and, on the moderate, station-specific datasets typical of optical observatories, are prone to shortcut learning: they may pick up background brightness, dome structures, or the timestamp printed on the image instead of the shape of the depletion. A lighter alternative is Two-Dimensional Principal Component Analysis (2DPCA) \cite{yang2004two}, which works directly on the image matrix and keeps part of its spatial structure. Yacoub et al. \cite{yacoub2025automatic} and Siddiqui \cite{siddiqui2025enhanced} recently combined 2DPCA with tree-based and ensemble classifiers for airglow images.

Three practical problems remain. First, local contrast enhancement turns dark cloud edges into sharp bands that classifiers confuse with EPBs, so cloudy frames have to be removed before classification. Second, consecutive exposures, taken $3\text{--}5$ minutes apart, are strongly correlated; random frame-level splits can therefore overestimate performance unless the model is also tested on held-out time windows and on stations not seen during training. Third, all-sky images carry the station code and timestamp burned into their corners, and one must check that the classifier does not rely on them.

This paper proposes a two-stage pipeline for EPB detection in $630.0\text{ nm}$ images from the Brazilian EMBRACE network. The first stage is a sky quality check, based on star counts and sky brightness, that discards cloud-covered frames. The second stage normalizes the background and local contrast, projects the image onto a small 2DPCA subspace, and classifies it with a stacking ensemble \cite{wolpert1992stacked}. We evaluate the pipeline with stratified, time-grouped, and text-masked cross-validation, with zero-shot transfer to an unseen station, and on full-night sequences, and we inspect its decisions with 2DPCA reconstruction residuals and SHapley Additive exPlanations (SHAP) \cite{lundberg2017unified}. Section~\ref{sec:methodology} describes the data and the method, Section~\ref{sec:results} presents the results, and Section~\ref{sec:conclusion} concludes the paper.
